\documentclass[12pt,a4paper]{article}
\usepackage[utf8]{inputenc}
\usepackage{graphicx}
\usepackage{float}
\usepackage{amsmath}

\title{Informe Taller 3: Búsqueda Anytime con AW-A*}
\author{Jose Salazar, Sofía Sepúlveda}
\date{\today}

\begin{document}
\maketitle

\section{Introducción}
El algoritmo \textit{Anytime Weighted A*} (AW-A*) es una extensión del clásico A* diseñada para escenarios en los que el tiempo de cómputo es limitado y se requiere obtener soluciones viables de manera inmediata. Su característica principal es la capacidad de entregar una primera solución rápidamente y luego refinarla progresivamente a medida que transcurre el tiempo disponible. Esto lo convierte en un enfoque especialmente útil en sistemas de tiempo real, como navegación robótica, logística dinámica o la toma de decisiones en videojuegos.

\section{Desarrollo Experimental}
Se realizaron pruebas en mapas de distinta escala y complejidad. En los escenarios pequeños, como los de $41 \times 41$, se observó que el algoritmo alcanza la solución óptima casi de inmediato, dado que el espacio de búsqueda es reducido. En consecuencia, el carácter \textit{anytime} no se aprecia con claridad, ya que la primera ruta encontrada corresponde prácticamente a la mejor solución posible.

En contraste, los experimentos en mapas de gran tamaño mostraron con mayor fuerza el valor de AW-A*. En mapas de $1500 \times 1300$ y $2000 \times 2000$, el algoritmo entregó soluciones iniciales aceptables y, con mayor tiempo, logró rutas de mejor calidad. La heurística Manhattan se comportó de manera más eficiente que la Euclidiana, puesto que permitió un mejor equilibrio entre calidad de la ruta y tiempo de cómputo. Además, valores de $w_{start}$ intermedios (por ejemplo, $2.5$) ofrecieron una buena combinación entre rapidez y capacidad de refinamiento posterior.

\begin{table}[H]
    \centering
    \resizebox{\textwidth}{!}{%
    \begin{tabular}{|c|c|c|c|c|c|c|c|c|}
    \hline
    \textbf{Size} & \textbf{Manhattan} & \textbf{w\_start} & \textbf{Solved Rate} & \textbf{Avg. BSL} & \textbf{Avg. BST} & \textbf{Avg. FST} & \textbf{Avg. FSL} \\
    \hline
    (1500,1300) & False & 2.0 & 0.5 & 12827.0 & 35.2 & 10.2 & 12875.0 \\
    (1500,1300) & False & 2.5 & 0.5 & 12843.0 & 40.9 &  9.6 & 12875.0 \\
    (1500,1300) & False & 3.0 & 0.5 & 12843.0 & 51.1 &  8.9 & 12875.0 \\
    (1500,1300) & True  & 2.0 & 0.5 & 12827.0 & 34.9 & 10.2 & 12875.0 \\
    (1500,1300) & True  & 2.5 & 0.5 & 12843.0 & 42.6 &  9.7 & 12875.0 \\
    (1500,1300) & True  & 3.0 & 0.5 & 12843.0 & 51.0 &  8.8 & 12875.0 \\
    (2000,2000) & False & 2.0 & 0.5 & 16739.7 & 52.1 & 22.8 & 16751.0 \\
    (2000,2000) & False & 2.5 & 0.5 & 16744.3 & 55.7 & 20.2 & 16751.0 \\
    (2000,2000) & False & 3.0 & 0.5 & 16796.3 & 59.9 & 16.6 & 16897.0 \\
    (2000,2000) & True  & 2.0 & 0.5 & 16739.7 & 52.1 & 22.8 & 16751.0 \\
    (2000,2000) & True  & 2.5 & 0.5 & 16744.3 & 55.3 & 20.2 & 16751.0 \\
    (2000,2000) & True  & 3.0 & 0.5 & 16796.3 & 59.5 & 16.4 & 16897.0 \\
    \hline
    \end{tabular}}
    \caption{Resultados en mapas grandes: tasa de éxito, longitudes promedio y métricas de tiempos.}
    \label{tab:mapas_grandes_full}
\end{table}

En lo referente a la comparación con otros algoritmos, AW-A* superó claramente a Greedy Best-First, que en nuestras ejecuciones no logró entregar soluciones válidas en mapas densos. A* por su parte alcanzó la optimalidad, pero con tiempos de cómputo prohibitivos. AW-A* se ubica en un punto intermedio, entregando rápidamente rutas utilizables y refinándolas con el tiempo, lo que lo convierte en una alternativa más práctica.

\section{Conclusiones}
Los experimentos confirman que AW-A* constituye un compromiso efectivo entre rapidez y calidad. Mientras en escenarios pequeños converge rápidamente a la solución óptima, en mapas de gran escala logra mantener una solución disponible en todo momento y mejorarla conforme se dispone de más tiempo. Sus limitaciones más relevantes son la sensibilidad a los parámetros de configuración y el costo creciente en mapas de gran tamaño.  

Entre sus aplicaciones reales destacan los sistemas de navegación autónoma, la logística de última milla y los videojuegos, donde contar con soluciones rápidas y mejorables resulta más valioso que esperar indefinidamente por la ruta perfecta.

\end{document}
